\chapter{Conclusion}
\label{sec:conclusion}

This thesis built and evaluated an agentic multimodal retrieval-augmented generation
system for questions about scientific papers. A planner fetches text, tables and
figures over several rounds; a Sufficiency Gate decides when the evidence suffices; a
Faithfulness Gate verifies every claim against its cited evidence; and every published
claim receives a proof box on the page and a place in an execution trace. The system
was evaluated on 150 SPIQA questions against a single-pass and a text-only baseline,
with a local 4-bit 7B model and, behind the identical graph, a stronger hosted model.

\paragraph{RQ1.} With the 7B model, the agentic loop with its gates was wrong on
\nFullWrong\% of the questions against \nSPWrong\% for single-pass, a reduction beyond
luck; but it was also correct less often (\nFullCorrect\% against \nSPCorrect\%), because
its abstentions did not single out the hard questions. On the small multi-hop subset it
was both less often wrong and more often right. With a stronger model behind the same
graph, abstention became selective.

\paragraph{RQ2.} Visual retrieval located the reference table or figure for 73\% of the
questions. Seeing tables and figures rather than their extracted text changed overall
accuracy little with the 7B model, but helped clearly on multi-hop questions.

\paragraph{RQ3.} When a claim cited the right table or figure, its proof box pointed at
the answer \nProofInsideFull\% of the time, and every checked table-cell proof did. Grounding
per question was limited by abstentions and by claims citing other sources, and IoU
proved a poor yardstick when proof and annotation are drawn at different granularity.

\section{Contributions}

\begin{itemize}
    \item An agentic multimodal RAG design with two gates and verified claim links,
    implemented end to end and run on consumer hardware with an open model.
    \item Exact proof boxes --- table cells, chart bars read from the PDF, quoted
    sentences, captions and diagram labels --- with the retriever's similarity map
    kept only for photographs.
    \item A three-way evaluation (correct, wrong, abstained) with a checked model judge,
    a selectivity analysis of abstention, and a grounding evaluation on hand-annotated
    answer regions.
    \item Evidence that the benefit of the gates depends on the capacity of the model
    behind them, obtained by holding the architecture fixed and changing the model.
\end{itemize}

\section{Future Work}

The results point to three directions. First, a hybrid in which the small local model
reads and writes while a stronger model makes the two judgements; it should be tested
on fresh questions from the unused SPIQA papers. Second, steering claims to cite the
table or figure that holds the answer rather than surrounding text, which would raise
grounding directly. Third, a user study measuring whether proof boxes and traces
actually help researchers verify answers, and an extension beyond computer-science
papers.
